\songtitle{No me vendas dolor de patria}

\tag*{2015}\tag{tropical-pop}\tag{spanish-lyrics}\tag{original-source}\tag{let-the-border-breathe-again}
\begin{notes}
Lyrics by Suno based on a Claude summary of Carlos Thompson's 2015 blog post \emph{Dolor de patria}. Spanish-language companion to \emph{Let the Border Breathe Again}, generated in the same session -- localizes the argument to the Colombia--Venezuela border, naming the Corinto eviction of a Nasa family and four Colombian presidents by name. Retrieved from Gabriel's Suno page (V6-MINI, 17 September 2026).
\end{notes}
\begin{style}
Tropical pop with candy-pop hooks, upbeat dembow-inspired groove, bright marimba plucks, clean electric guitar, bubbly synth bass, handclaps, crisp modern radio production; conversational Spanish lead vocal with playful verses, tight melodic chorus, brief half-time bridge, energetic final refrain
\end{style}
\begin{lyrics}
\songpart{Intro}
No todo grito es verdad,\\
no toda bandera sabe cuidar.

\songpart{Verse 1}
Primero, pon los nombres sobre la mesa,\\
no todo colombiano es guerra ni amenaza.\\
Hay grupos ilegales, sí, cruzando la frontera,\\
pero un sello en un discurso no explica la historia entera.

Segundo, una expulsión puede tener ley,\\
los gobiernos deportan, aquí también lo sé.\\
En Corinto cayó una casa nasa aquel año,\\
y el silencio pesó más que cualquier titular extraño.

\songpart{Pre-Chorus}
Pero entre tratados viejos y cédulas dudosas,\\
la vida no cabe en frases luminosas.\\
Libre tránsito, voto comprado, papeles sin final,\\
la frontera es una herida con memoria legal.

\songpart{Chorus}
No me vendas dolor de patria,\\
si eliges qué dolor mirar.\\
Tres mil dentro de la casa,\\
¿quién los vino a reclamar?\\
No me vendas orgullo herido,\\
con la rabia en oferta especial.\\
Si la bandera tapa al vecino,\\
¿de qué sirve gritar nacional?

\songpart{Verse 2}
País A es pobre, país B tiene más,\\
X busca trabajo, Y teme perder su lugar.\\
Subsidios, gasolina, mercado y contrabando,\\
el que llega se vuelve invasor en el relato.

Pero X también puede huir de una finca,\\
de una guerrilla, de una amenaza clandestina.\\
No todos cruzan para ganar ventaja,\\
algunos cruzan porque quedarse era la trampa.

\songpart{Pre-Chorus}
Y Maduro juega fuerte con la cámara encendida,\\
estado de excepción, redada dirigida.\\
Antes de las elecciones, fabrica un enemigo ideal,\\
la frontera convertida en escenario electoral.

\songpart{Chorus}
No me vendas dolor de patria,\\
si eliges qué dolor mirar.\\
Tres mil dentro de la casa,\\
¿quién los vino a reclamar?\\
No me vendas orgullo herido,\\
con la rabia en oferta especial.\\
Si la bandera tapa al vecino,\\
¿de qué sirve gritar nacional?

\songpart{Bridge}
Uribe, Pastrana, Samper, Santos,\\
cuántos años dejando el problema intacto.\\
El desplazado sin cámara, el migrante sin voz,\\
la misma deuda cambiando de color.

Santos y Holguín: hablen, negocien,\\
que ningún soldado la frontera desbroce.\\
Ni renuncia teatral ni desfile marcial,\\
la salida sensata también puede ser bilateral.

\songpart{Final Chorus}
No me vendas dolor de patria,\\
ni patriotero sentimental.\\
Mira al que pierde dentro de casa,\\
antes de posar para el altar.\\
No me vendas orgullo herido,\\
con titulares para incendiar.\\
Si la bandera tapa al vecino,\\
yo no la voy a levantar.

\songpart{Outro}
No todo grito es verdad,\\
no toda bandera sabe cuidar.\\
Primero la gente, después el papel,\\
la patria se demuestra cuando sabe responder.
\end{lyrics}

